

\begin{frame}[fragile]
\frametitle{Rappels: Notions et vocabulaire}

Vous devriez savoir: 
\begin{redbullet}
  \item \alert{Item (entrée)}, \alert{document} en ce qui nous concerne.
  
  \item \alert{Fonction de Map}, implante la transformation appliquée à un item pendant la phase de Map.

  \item \alert{Paire intermédiaire}, c'est le résultat de la fonction de Map. 
  
  Une paire intermédiaire $(k, v)$ comprend l'identifiant (étiquette) du groupe auquel
  $v$ appartient.

  \item \alert{Groupe intermédiaire}, ensemble des valeurs produites par la fonction
  de Map et partageant le même identifiant de groupe.
  
    \item \alert{Fonction de Reduce}, implante la transformation appliquée un groupe
    pendant la phase de Reduce.
    
\end{redbullet}

\alert{Ne pas retenir (pour l'instant)}: comment tout cela s'exécute en distribué, avec
parallélisation et reprise sur panne.

\end{frame}


\section{Frameworks MapReduce: MongoDB}



\begin{frame}[fragile]
\frametitle{La notion de \textit{framework}}

\textit{Framework} = environnement assumant les aspects \alert{génériques} d'un traitement.

\vfill
Ici: tout ce qui n'est pas le cuisinier, le pressoir .

\figSlideWithSize{jusfruits+}{9cm}

\begin{redbullet}
  \item Impose un modèle d'exécution contraint.
  \item Mais il suffit de fournir la spécification \alert{fonctionnelle}: $F_{map}()$ et $F_{red}()$
\end{redbullet}

\end{frame}


\begin{frame}[fragile]
\frametitle{Exemple: compter les fruits sains, par catégorie}

$F_{map}()$ qui teste si un produit 
est pourri et \alert{émet} une paire intermédiaire.

\begin{minted}[baselinestretch=.9,
               fontsize=\small,
               framesep=2mm]{javascript}
function controleFruit (fruit)
{
   if (fruit.statut != pourri) {
     emit (fruit.type, 1)
   }
}
\end{minted}


Chaque groupe intermédiaire contient autant de 1 que de fruits sains.
\vfill

$F_{red}()$  additionne
les valeurs dans un groupe.

\begin{minted}[baselinestretch=.9,
               fontsize=\small,
               framesep=2mm]{javascript}
function compterPomme (typeFruit, groupe)
{
  return <typeFruit, sum(groupe)>
}
\end{minted}

La fonction est appliquée à \alert{chaque} groupe.

\smallskip

Il reste à soumettre ces deux fonctions au \textit{framework}. \alert{Clair?}

\end{frame}

\section{Application avec MongoDB}


\begin{frame}[fragile]
\frametitle{MapReduce et MongoDB}

MongoDB comprend un moteur de calcul MapReduce. 

\vfill


Les fonctions doivent être écrites en Javascript.

\vfill

Dans ce qui suit: deux exemples de traitements MapReduce pour

\begin{redbullet}
  \item Comprendre les principes, le raisonnement.
  \item Réaliser les limites quand il s'agit d'obtenir des traitements un peu complexes.
  \item Méditer sur l'utilité de la chose...
\end{redbullet}

Fonctions fournies: à tester avec l'utilitaire \texttt{mongo} ou avec RoboMongo (ou tout
autre client).
\end{frame}

\begin{frame}[fragile]
\frametitle{Exemple: regroupement de films}

On veut regrouper les films par réalisateur

\vfill


La fonction $F_{map}$.

\begin{minted}[baselinestretch=.9,
               fontsize=\small,
               framesep=2mm]{javascript}
var mapRealisateur = function() {
                  emit(this.director._id, this.title);
          };
\end{minted}

\vfill

La fonction $F_{red}$.

\begin{minted}[baselinestretch=.9,
               fontsize=\small,
               framesep=2mm]{javascript}
   var reduceRealisateur = function(directorId, titres) {
      var res = new Object();
      res.director = directorId;
      res.films = titres;
      return res;
    };
\end{minted}

\vfill

Soumission au \textit{framework}

\begin{minted}[baselinestretch=.9,
               fontsize=\small,
               framesep=2mm]{javascript}
 db.movies.mapReduce(mapRealisateur,  reduceRealisateur, {out: {"inline": 1}} )
\end{minted}
\end{frame}

\begin{frame}[fragile]
\frametitle{Le résultat}

Chaque exécution de $F_{red}$ produit un document de la forme.
\begin{minted}[baselinestretch=.9,
               fontsize=\small,
               framesep=2mm]{javascript}
{
	"_id" : "artist:3",
	"value" : {
		"director" : "artist:3",
		"films" : [
			"Vertigo",
			"Psychose",
			"Les oiseaux",
			"Pas de printemps pour Marnie",
			"La mort aux trousses"
		]
	}
}
\end{minted}

\vfill

On aimerait connaître le réalisateur: exercice!

\vfill

Options d'exécution: voir le support.


\end{frame}


\section{Jointures avec MapReduce}


\begin{frame}[fragile]
\frametitle{Jointure: la méthode}

Notre
base est celle contenant les films \alert{avec références}. 

\vfill

Films et artistes sont dans des documents \alert{distincts}.
Le but est  d'obtenir pour chaque artiste la liste des films qu'il/elle a réalisé.

\vfill

\alert{Principe (général)}:  on exploite 
le mécanisme de regroupement pour associer, dans un même groupe les informations à combiner.

\vfill

On va créer autant de groupes que d'artiste. Dans chaque groupe on place:

\begin{redbullet}
  \item l'artiste dont l'identifiant correspond à l'identifiant du groupe;
  \item les films (0, 1 ou plusieurs) dont l'identifiant \alert{du metteur en scène}
    correspond à l'identifiant du groupe.
\end{redbullet}

Méthode très représentative de l'application de  MapReduce à des algorithmes complexes (pas trop quand même).

\end{frame}


\begin{frame}[fragile]
\frametitle{Illustration: groupement d'un metteur en scène et de ses films}

\figSlideWithSize{jointure-mr}{10cm}
\end{frame}

\begin{frame}[fragile]
\frametitle{En pratique: la fonction de Map}

\begin{minted}[baselinestretch=.9,
               fontsize=\small,
               framesep=2mm]{javascript}
   var mapJoin = function() {
      // Est-ce que l'id du document contient le mot "artist"?
      if (this._id.indexOf("artist") != -1) {
        // Oui ! C'est un artiste. Ajoutons-lui son type.
        this.type="artist";        
        // On produit une paire avec pour id celle de l'artiste
        emit(this._id, this);
      }
      else {
        // Non: c'est un film. Ajoutons-lui son type.
       this.type="film";
       // Simplifions un peu le document pour l'affichage
       delete this.summary;
       delete this.actors;
        // On produit une paire avec pour id celle du metteur en sc.
       emit(this.director._id, this);
     }
    };
\end{minted}

\end{frame}


\begin{frame}[fragile]
\frametitle{En pratique: la fonction de Reduce}

\begin{minted}[baselinestretch=.9,
               fontsize=\small,
               framesep=2mm]{javascript}
   var reduceJoin = function(id, items) {
   
     var director = null, films={result: []}
      
     // On cherche l'artiste dans cette liste    
     for (var idx = 0; idx < items.length; idx++) {
       if (items[idx].type=="artist") {
            director = items[idx];
        }
     }             
           
     // Maintenant, 'director' contient l'artiste : on l'affecte aux films   
     for (var idx = 0; idx < items.length; idx++) {
        if (items[idx].type=="film"  && director != null) {
            items[idx].director = director;
            films.result.push (items[idx]);
         }             
      }
      return films; 
    };
\end{minted}

\end{frame}